\documentclass[11pt]{article}

\usepackage[letterpaper,margin=0.92in]{geometry}
\usepackage[T1]{fontenc}
\usepackage[utf8]{inputenc}
\usepackage{booktabs}
\usepackage{enumitem}
\usepackage{fancyhdr}
\usepackage{microtype}
\usepackage{setspace}

\setstretch{1.05}
\setlength{\parindent}{0pt}
\setlength{\parskip}{5pt}
\setlist[enumerate]{leftmargin=2.2em,itemsep=3pt,topsep=3pt}
\setlist[itemize]{leftmargin=2.0em,itemsep=2pt,topsep=3pt}

\pagestyle{fancy}
\fancyhf{}
\fancyhead[C]{\footnotesize\textsc{Defensive Lineup Legality and Continuing Compliance Code}}
\fancyfoot[C]{\thepage}
\renewcommand{\headrulewidth}{0.4pt}
\renewcommand{\footrulewidth}{0pt}

\newcommand{\articleheading}[1]{%
  \vspace{6pt}
  \begin{center}
    \large\bfseries #1
  \end{center}
  \vspace{-3pt}
}

\newcommand{\sectionheading}[2]{%
  \vspace{4pt}
  \noindent\textbf{\S\,#1.\quad #2}\par
}

\begin{document}

\thispagestyle{empty}
\begin{center}
{\small DOCUMENT NO.\ SFL--LCC--7I/2026}\\[0.35in]

{\Large\bfseries THE CO-ED SLOWPITCH SOFTBALL}\\[4pt]
{\Large\bfseries DEFENSIVE LINEUP LEGALITY, COMPOSITION,}\\[4pt]
{\Large\bfseries ELIGIBILITY, AND CONTINUING COMPLIANCE CODE}\\[0.18in]

{\large (the ``Lineup Code'')}\\[0.35in]

\rule{0.72\textwidth}{0.6pt}\\[0.16in]
{\small\textsc{Effective July 2026}}\\[0.16in]
\rule{0.72\textwidth}{0.6pt}
\end{center}

\vspace{0.35in}

\textbf{PREAMBLE AND RECITALS}

\textbf{WHEREAS}, the orderly administration of co-ed slowpitch softball
requires that each defensive half-inning be contested by a properly
constituted, positionally complete, gender-compliant, and duly authorized
fielding complement;

\textbf{WHEREAS}, the unauthorized occupation, abandonment, duplication, or
improvised reclassification of any defensive position threatens the integrity
of the lineup as a whole and shall not be excused merely because the ball has
already been put into play;

\textbf{WHEREAS}, compliance is not a matter of substantial approximation,
customary practice, good intentions, optimistic head-counting, or the
manager's sincere belief that the arrangement ``is probably close enough'';

\textbf{WHEREAS}, the requirements stated herein are cumulative, inning
specific, and of continuing force, such that compliance during one inning
creates no presumption whatsoever of compliance during any subsequent inning;
and

\textbf{WHEREAS}, failure to observe these requirements may subject the
offending lineup, its preparers, and the Club on whose behalf it is tendered
to immediate corrective action and such further consequences as the
competent authority may deem necessary, appropriate, exemplary, or
uncomfortably memorable;

\textbf{NOW, THEREFORE}, the following provisions are hereby declared to
govern the preparation, review, issuance, and continued use of every
seven-inning defensive lineup.

\articleheading{ARTICLE I\\DEFINITIONS AND RULES OF CONSTRUCTION}

\sectionheading{1.01}{Defined Terms.}
For purposes of this Lineup Code, the following terms shall have the meanings
set forth below, whether used in the singular or plural:

\begin{enumerate}[label=(\alph*)]
  \item \textbf{``Active Position''} means a defensive position that is
  required to be occupied under the applicable Fielding Complement.

  \item \textbf{``Available Player''} means a rostered player affirmatively
  designated as available for the game in question. Mere presence in the
  vicinity of the field shall not, standing alone, constitute availability.

  \item \textbf{``Defensive Inning''} means each of the seven separately
  regulated innings for which a defensive assignment is issued.

  \item \textbf{``Fielding Complement''} means the legally permissible group
  of eight, nine, or ten players assigned to the defense in a Defensive
  Inning.

  \item \textbf{``Infield''} means P, C, 1B, 2B, 3B, and SS. For the avoidance
  of doubt, and notwithstanding any contrary dugout folklore, both P and C
  are Infield positions.

  \item \textbf{``Outfield''} means LF, LC, RC, and RF.

  \item \textbf{``Position Preference''} means a position for which a player
  has been affirmatively declared eligible. Position Preferences are
  unranked and equally preferred, but they are not optional suggestions.

  \item \textbf{``Woman Player''} means an Available Player identified as a
  woman in the roster data governing the applicable game.
\end{enumerate}

\sectionheading{1.02}{Strict and Cumulative Construction.}
Every ``shall,'' ``must,'' and ``may not'' appearing herein is mandatory.
Each requirement applies in addition to, and not in substitution for, every
other requirement. Satisfaction of the numerical requirement shall not cure a
geographic defect; satisfaction of the geographic requirement shall not cure
an eligibility defect; and no aggregation of partial compliance shall be
recognized as the legal equivalent of complete compliance.

\articleheading{ARTICLE II\\AUTHORIZED FIELDING COMPLEMENTS}

\sectionheading{2.01}{Ten-Player Complement.}
A ten-player Fielding Complement shall occupy all of the following positions:
P, C, 1B, 2B, 3B, SS, LF, LC, RC, and RF. It shall include no fewer than four
Woman Players, of whom no fewer than one shall be stationed in the Infield and
no fewer than one shall be stationed in the Outfield.

\sectionheading{2.02}{Nine-Player Complement.}
If a lawful ten-player Fielding Complement cannot be constituted because
fewer than ten players are available or because fewer than four Woman Players
are available, a nine-player Fielding Complement may be constituted only if:
\begin{enumerate}[label=(\roman*)]
  \item at least nine players are available;
  \item at least three of those players are Woman Players;
  \item RF is left wholly and conspicuously unoccupied;
  \item each of P, C, 1B, 2B, 3B, SS, LF, LC, and RC is occupied exactly once;
  and
  \item at least one Woman Player occupies the Infield and at least one Woman
  Player occupies the Outfield.
\end{enumerate}
No tenth player may be introduced as a rover, auxiliary fielder, informal
short-fielder, or person who is ``just standing over there.''

\sectionheading{2.03}{Eight-Player Complement.}
An eight-player Fielding Complement may be constituted only if exactly eight
players, including at least three Woman Players, are available for assignment.
In such event, C and RF shall remain unoccupied; each of P, 1B, 2B, 3B, SS,
LF, LC, and RC shall be occupied exactly once; at least one Woman Player shall
occupy the Infield; and at least one Woman Player shall occupy the Outfield.

\sectionheading{2.04}{Mandatory Maximum Lawful Complement.}
The largest lawful Fielding Complement supported by the available roster shall
govern. A smaller complement may not be selected merely for convenience,
aesthetic symmetry, strategic experimentation, avoidance of difficult
decisions, or relief from the administrative burden of counting.

\sectionheading{2.05}{Absolute Minimum.}
Fewer than eight Available Players, or fewer than three available Woman
Players, is legally insufficient to constitute any Fielding Complement
recognized under this Lineup Code. No amount of enthusiasm, positional
versatility, arithmetic creativity, or informal consensus shall cure such
deficiency.

\articleheading{ARTICLE III\\PER-INNING ASSIGNMENT REQUIREMENTS}

\sectionheading{3.01}{Exact Occupancy.}
For each Defensive Inning, every Active Position shall be assigned to exactly
one eligible Available Player. No Active Position may be vacant, shared,
alternated within the same assignment, reserved for later determination, or
filled by a player identified only as ``TBD.''

\sectionheading{3.02}{Single-Position Limitation.}
No player may occupy, be credited with, or be assigned responsibility for more
than one defensive position during the same Defensive Inning. Athletic range,
confidence, prior experience, and an expressed willingness to ``cover both''
shall have no legal effect.

\sectionheading{3.03}{Inactive Positions.}
Every position rendered inactive by Article II shall remain empty for the
entire applicable Defensive Inning. The omission of C or RF is a prescribed
structural condition, not an invitation to rename, relocate, or disguise an
additional fielder.

\sectionheading{3.04}{Independent Inning Review.}
The requirements of this Article shall be independently satisfied in each of
the seven Defensive Innings. A player substitution, positional transfer, or
return from sitting out shall trigger renewed compliance review. Legality does
not vest permanently upon acceptance of the first-inning lineup.

\articleheading{ARTICLE IV\\POSITIONAL ELIGIBILITY}

\sectionheading{4.01}{Preference as Authorization.}
Except as expressly provided in \S\,4.02, a player may be assigned only to a
position contained in that player's recorded Position Preferences. Silence,
blank data, assumed competence, necessity, and the manager's personal
assessment that a player ``can probably handle it'' shall not constitute
authorization.

\sectionheading{4.02}{Limited RF-to-RC Operation of Law.}
Solely when an eight- or nine-player Fielding Complement is in effect, a
recorded preference for RF shall automatically confer eligibility for RC for
that game. This narrowly drawn legal fiction exists because RF is inactive in
both such complements.

\sectionheading{4.03}{No Implied Expansion.}
The rule in \S\,4.02 shall not be construed to:
\begin{enumerate}[label=(\alph*)]
  \item confer RF eligibility from an RC preference;
  \item confer eligibility for LC, LF, or any Infield position;
  \item operate in a ten-player Fielding Complement; or
  \item establish any broader doctrine of ``close enough'' positional
  equivalence.
\end{enumerate}

\articleheading{ARTICLE V\\PRIORITY, ADMINISTRATION, AND RECORDS}

\sectionheading{5.01}{Supremacy of Legality.}
The requirements of Articles II through IV are conditions precedent to a valid
lineup. Goals concerning equal playing time, continuous assignments, player
satisfaction, or limiting the number of distinct positions are subordinate
and shall not excuse, offset, or retroactively sanitize an unlawful
assignment.

\sectionheading{5.02}{Roster and Availability Record.}
Before a lineup is generated, the responsible manager shall maintain an
accurate record of player names, gender designations, Position Preferences,
and game-specific availability. Any lineup derived from materially inaccurate
input shall be treated as suspect from inception and shall remain subject to
immediate review, reformation, and replacement.

\sectionheading{5.03}{Burden of Verification.}
The person presenting or using a lineup bears the continuing burden of
verifying its compliance. Reliance upon software, a prior lineup, a
spreadsheet, an unreviewed text message, or the confident assurances of
another person shall not transfer, diminish, or extinguish that burden.

\sectionheading{5.04}{No Waiver by Inattention.}
The failure of any player, manager, spectator, opposing participant, or other
interested person to identify a defect immediately shall not constitute
approval, ratification, waiver, estoppel, acquiescence, or evidence that the
defect was too minor to matter.

\articleheading{ARTICLE VI\\NONCOMPLIANCE AND CONSEQUENCES}

\sectionheading{6.01}{Material Breach.}
Any violation of Articles II, III, or IV shall constitute a material lineup
defect. A defective assignment is noncompliant \emph{ab initio} and shall not
be rehabilitated by competent play, favorable results, lack of objection, or
the subsequent discovery that no batted ball reached the improperly arranged
portion of the field.

\sectionheading{6.02}{Corrective Authority.}
Upon actual or suspected noncompliance, the competent authority may require
immediate cessation, compulsory reassignment, submission of a replacement
lineup, removal of an ineligible assignment, formal explanation, renewed
certification, or any combination thereof. Nothing herein shall require that
advance notice or a leisurely opportunity to cure be afforded before such
measures are imposed.

\sectionheading{6.03}{Reservation of Consequences.}
The remedies stated in \S\,6.02 are not exclusive. A violation may expose the
lineup and the Club to disallowance, nullification, forfeiture, adverse
determination, reputational censure, and such additional sanctions,
consequences, or solemn administrative disappointment as may lie within the
jurisdiction of the governing authority. The absence of a separately
enumerated penalty shall not be interpreted as permission.

\articleheading{ARTICLE VII\\FINAL PROVISIONS}

\sectionheading{7.01}{Severability.}
If any provision of this Lineup Code is determined to be invalid or
unenforceable, the remaining provisions shall continue in full force, and the
invalid provision shall be construed as narrowly as necessary to preserve the
maximum lawful degree of lineup regulation.

\sectionheading{7.02}{Amendment.}
No amendment, exception, side agreement, dugout proclamation, or purported
waiver shall be effective unless incorporated into the governing lineup rules
and, where applicable, the optimization system itself. Oral modifications
made while the teams are taking the field are strongly disfavored.

\sectionheading{7.03}{Headings.}
Article and section headings are supplied for organizational convenience and
shall not be invoked to evade the plain, stern, and unmistakable meaning of
the operative text.

\sectionheading{7.04}{Effective Scope.}
This Lineup Code applies to one game at a time and to all seven Defensive
Innings of each such game. It shall take effect upon preparation of the
lineup and remain in force until the final defensive assignment has concluded
or all compliance questions have been resolved, whichever occurs later.

\vfill
\begin{center}
\rule{0.82\textwidth}{0.5pt}\\[8pt]
\textbf{CERTIFICATION OF LINEUP COMPLIANCE}\\[6pt]
\end{center}

By presenting a defensive lineup for use, the undersigned certifies, under the
full weight of managerial responsibility, that the lineup has been reviewed
for exact positional occupancy, player eligibility, minimum participation by
Woman Players, required Infield and Outfield placement, and all other matters
set forth herein.

\vspace{0.32in}
\begin{tabular}{@{}p{0.45\textwidth}@{\hspace{0.08\textwidth}}p{0.37\textwidth}@{}}
\hrulefill & \hrulefill\\[-2pt]
Responsible Manager & Date and Time of Certification
\end{tabular}

\end{document}
